\documentclass[11pt]{article}
\usepackage[margin=1in]{geometry}
\usepackage{amsmath}
\usepackage{booktabs}
\usepackage{hyperref}
\usepackage{graphicx}
\usepackage{natbib}
\usepackage{xcolor}

\title{Deciding When a Group-Chat Bot Should Speak:\\
A Reproducible Gate for Addressedness}
\author{Weijun Sun\\
Central South University, Changsha, China\\
\href{https://github.com/Weijun-Sun08/qq-addressedness-gate}{\texttt{Weijun-Sun08/qq-addressedness-gate}}}
\date{2026}

\begin{document}
\maketitle

\begin{abstract}
Group-chat agents fail in a characteristic way: they answer messages that were never addressed to them.
We study the addressedness decision in isolation---given a message, its speaker, and a short context, should a bot reply?---and release a small, dependency-free implementation with three replaceable parts: a rule layer, an HTTP judge that answers the binary question, and an OpenAI-compatible chat client.
On 574 human-confirmed messages from a deployed assistant, a rules-only policy reaches 76.8\% accuracy with high precision (94.1\%) but low recall (50.4\%), while a 7B judge reaches 96.5\% accuracy, 99.2\% precision and 92.9\% recall at a median latency of 0.46\,s on CPU.
A character n-gram classifier matches the judge in-distribution (97.6\%) but drops to 87.8\% under a temporal split, whereas the judge's errors concentrate on third-person mentions and image messages.
We release the evaluation harness, the metric definitions, and the dataset schema so that addressedness strategies can be compared without rebuilding the surrounding bot.
\end{abstract}

\section{Introduction}
Conversational agents deployed in group chats must decide, for every message, whether they are being addressed.
Answering too often is intrusive; answering too rarely makes the agent useless.
Most frameworks fold this decision into the transport and the generator, which makes strategies hard to compare and results hard to reproduce.

We isolate the decision as a binary classification problem and provide a reference implementation with a deliberately narrow interface.
Our contributions are:
(i) a formulation of addressedness gating with explicit failure semantics (silence when the judge is unavailable);
(ii) a small library whose judge and generator are interchangeable; and
(iii) an offline evaluation on real group-chat traffic comparing rules, a classical classifier, and a small language model.

\section{Related Work}
Multi-party dialogue requires tracking who is being addressed; addressee identification has been studied as a task in its own right.
Open-source frameworks such as Rasa~\citep{rasa2017} support building dialogue agents but do not expose the addressedness decision as a reusable component.
Deployed bots typically rely on mentions or keyword matching, which are cheap but brittle for third-person references.

\section{Method}
The gate has three layers.

\paragraph{Rules.}
A message is \emph{skipped} if it contains no alphanumeric or CJK characters (punctuation-only messages), and treated as an explicit call if it mentions the bot or contains a configured name.
All remaining messages are delegated to the judge.

\paragraph{Judge.}
The judge is an HTTP service: \texttt{POST /judge} with \texttt{\{who, text, ctx\}} returns \texttt{\{"answer": "yes"|"no"\}}.
The reference implementation serves a Qwen2.5-7B~\citep{qwen25} model adapted with LoRA~\citep{lora} and exported to OpenVINO INT4 on CPU.
If the service is unavailable, the gate stays silent.

\paragraph{Generation.}
When the gate decides to reply, the message is sent to an OpenAI-compatible \texttt{/chat/completions} endpoint, so llama.cpp, Ollama, vLLM, or a hosted API can be substituted without touching the gate.

\section{Evaluation}
\paragraph{Data.}
574 messages from a deployed QQ assistant, labelled by the operator, collected after the judge's training cutoff.
Positive (should reply) 252, negative 322.

\paragraph{Results.}
Table~\ref{tab:main} reports accuracy, precision, recall, and F1.
Rules alone are precise but blind: they recover only half of the positives.
Delegating ambiguous messages to the judge recovers recall without sacrificing precision; letting rules answer on a name match costs precision, because third-person mentions look like calls.

\begin{table}[t]
\centering
\caption{Addressedness decisions on 574 held-out messages. Rules only treat unclear cases as ``no''.}
\label{tab:main}
\begin{tabular}{lcccc}
\toprule
Strategy & Acc. & Prec. & Rec. & F1\\
\midrule
Rules only & 76.8 & 94.1 & 50.4 & 65.6\\
Judge only & 96.5 & 99.2 & 92.9 & 95.9\\
Rules (noise skip) + judge & 96.5 & 99.2 & 92.9 & 95.9\\
Rules (name first) + judge & 95.5 & 96.3 & 93.3 & 94.8\\
\bottomrule
\end{tabular}
\end{table}

\paragraph{Classifier baseline and temporal drift.}
A character n-gram TF-IDF classifier with logistic regression reaches 97.6\% in-distribution.
Trained on data before 2026-11-07 and tested after, it falls to 87.8\%, close to the majority-class rate of 71.4\% on that split, indicating reliance on surface phrasing rather than context.

\paragraph{Latency.}
Judge calls over HTTP take 463\,ms at p50, 1.3\,s at p90, and 2.9\,s at p99 on CPU.

\paragraph{Qualitative errors.}
The judge's two false positives are name-like mentions not directed at the bot; its eighteen false negatives are mostly third-person references and messages accompanying images.

\section{Discussion}
The results suggest three practical guidelines.
First, rules are better as a cheap \emph{skip} filter than as a decision maker.
Second, a small local model is sufficient for the decision; no large model is needed at this stage.
Third, reported accuracy is sensitive to prompt format: the same model reaches 98.6\% when fed its training template in-process, versus 96.5\% through the service contract, so papers should state which interface was evaluated.

\section{Ethics and Data}
The public repository contains no chat logs.
Evaluation data is supplied locally and never uploaded by the code.
The deployment that produced the evaluation data involved a private group; messages were anonymised before analysis.

\section{Availability}
Code, tests, and the evaluation harness are available at
\url{https://github.com/Weijun-Sun08/qq-addressedness-gate} under the MIT license.

\bibliographystyle{plainnat}
\bibliography{refs}

\end{document}

